\section{Block Lyapunov Identities for the Calibrated Clamp}
\label{app:block_identities}

This appendix isolates the algebra used in the paper's shorting argument.  The
only inputs are the block-triangular calibrated operator family, the trace-class
stationary covariance, and the Cameron--Martin finite-energy assumption on the
outgoing drift channel.  All covariance computations below are made for centered
variables; deterministic means play no role in the operator identities.

\subsection{Block identities and stationary convolutions}

Write $\mathcal{H}=\mathcal{H}_k\oplus\mathcal{H}_I$ with
$\mathcal{H}_k=\mathbb{R}\phi_k$.  Under the calibrated clamp,
\begin{equation}
 \mathcal{A}^{(\lambda)}=
 \begin{pmatrix}-\lambda&0\\a_{Ik}&\mathcal{A}_{II}\end{pmatrix},
 \qquad
 \mathcal{D}^{(\lambda)}=
 \begin{pmatrix}\varepsilon_0/\lambda&0\\0&D_{II}\end{pmatrix}.
 \label{eq:exact_blocks}
\end{equation}
Let $\Sigma_\lambda$ be the stationary covariance solving
\[
  \mathcal{A}^{(\lambda)}\Sigma_\lambda
  +\Sigma_\lambda(\mathcal{A}^{(\lambda)})^*
  +\mathcal{D}^{(\lambda)}=0.
\]
The background bulk covariance is
\begin{equation}
  Q_0
  =
  \int_0^\infty e^{r\mathcal{A}_{II}}D_{II}e^{r\mathcal{A}_{II}^*}\,dr,
  \label{eq:Q0_convolution}
\end{equation}
which is the stationary covariance of the decoupled bulk equation.

\begin{lemma}[Bulk domination by independent stationary convolutions]
\label{lem:bulk_domination_convolution}
The stationary bulk coordinate admits the decomposition
\begin{equation}
  X_I(t)=X_I^{(0)}(t)+Y_{\lambda,I}(t),
  \qquad
  Y_{\lambda,I}(t)
  =
  \int_{-\infty}^t
  e^{(t-s)\mathcal{A}_{II}}a_{Ik}X_k(s)\,ds,
  \label{eq:bulk_stationary_convolution}
\end{equation}
where
\begin{equation}
  X_I^{(0)}(t)
  =
  \int_{-\infty}^t
  e^{(t-s)\mathcal{A}_{II}}D_{II}^{1/2}\,dW_I(s).
  \label{eq:background_stationary_convolution}
\end{equation}
The two summands $X_I^{(0)}(t)$ and $Y_{\lambda,I}(t)$ are driven by independent
noises and hence have zero cross covariance.  Therefore
\begin{equation}
  \Sigma_{\lambda,II}
  =Q_0+\operatorname{Cov}(Y_{\lambda,I})\succeq Q_0.
  \label{eq:bulk_domination}
\end{equation}
\end{lemma}

\begin{proof}
The target coordinate is the stationary Ornstein--Uhlenbeck convolution
\begin{equation}
  X_k(t)
  =
  \int_{-\infty}^t e^{-\lambda(t-s)}
  \sqrt{\varepsilon_0/\lambda}\,dW_k(s),
  \label{eq:target_stationary_convolution}
\end{equation}
with $W_k$ independent of the bulk cylindrical Brownian motion $W_I$.  Variation
of constants for the lower triangular system gives
\eqref{eq:bulk_stationary_convolution}--\eqref{eq:background_stationary_convolution}.
The exponential stability assumptions guarantee convergence of the improper
stochastic and deterministic convolutions in $L^2(\Omega;\mathcal{H}_I)$.

For any $h,g\in\mathcal{H}_I$, the scalar random variable
$\langle h,X_I^{(0)}(t)\rangle$ is measurable with respect to the sigma-field
generated by $W_I$, whereas $\langle g,Y_{\lambda,I}(t)\rangle$ is measurable
with respect to the sigma-field generated by $W_k$.  These sigma-fields are
independent, and the centered variables therefore have zero covariance.  Hence
\[
  \operatorname{Cov}(X_I(t))
  =\operatorname{Cov}(X_I^{(0)}(t))
   +\operatorname{Cov}(Y_{\lambda,I}(t)).
\]
The first covariance is exactly \eqref{eq:Q0_convolution}.  The second is a
positive covariance operator, proving the Loewner domination
\eqref{eq:bulk_domination}.
\end{proof}

\begin{lemma}[Cross-covariance resolvent identity]
\label{lem:block_lyapunov_identities}
For every $\lambda>0$,
\begin{equation}
  \Sigma_{\lambda,kk}=\frac{\varepsilon_0}{2\lambda^2},
  \label{eq:Sigmakk_exact}
\end{equation}
and the target-to-bulk cross covariance, identified with the vector
$\Pi_I\Sigma_\lambda\phi_k\in\mathcal{H}_I$, is
\begin{equation}
  \Sigma_{\lambda,Ik}
  =\frac{\varepsilon_0}{2\lambda^2}
    (\lambda I-\mathcal{A}_{II})^{-1}a_{Ik}.
  \label{eq:SigmaIk_exact_vector}
\end{equation}
The resolvent in \eqref{eq:SigmaIk_exact_vector} exists both on
$\mathcal{H}_I$ and, under Assumption~\ref{assum:cm_stability}, on
$H_{Q_0}$.  Consequently $\Sigma_{\lambda,Ik}\in H_{Q_0}$.
\end{lemma}

\begin{proof}
The scalar target Lyapunov equation gives
\[
  -2\lambda\Sigma_{\lambda,kk}+\frac{\varepsilon_0}{\lambda}=0,
\]
which proves \eqref{eq:Sigmakk_exact}.  Since $\mathcal{A}_{II}$ generates an
exponentially stable semigroup, every $\lambda>0$ belongs to the resolvent set
of $\mathcal{A}_{II}$ and
\begin{equation}
  (\lambda I-\mathcal{A}_{II})^{-1}x
  =
  \int_0^\infty e^{-\lambda r}e^{r\mathcal{A}_{II}}x\,dr,
  \qquad x\in\mathcal{H}_I.
  \label{eq:resolvent_laplace_HI}
\end{equation}
The same formula holds on $H_{Q_0}$ for the restricted exponentially stable
semigroup from Assumption~\ref{assum:cm_stability}.  In particular,
$(\lambda I-\mathcal{A}_{II})^{-1}a_{Ik}$ is a well-defined element of
$H_{Q_0}$.

It remains to compute the sign and the scalar coefficient.  By
Lemma~\ref{lem:bulk_domination_convolution}, the background bulk convolution is
independent of $X_k$, so only $Y_{\lambda,I}$ contributes to the cross covariance.
Using stationarity of the target Ornstein--Uhlenbeck process,
\[
  \mathbb{E}\bigl[X_k(t-r)X_k(t)\bigr]
  =\Sigma_{\lambda,kk}e^{-\lambda r},
  \qquad r\ge0.
\]
Therefore
\begin{align*}
  \Sigma_{\lambda,Ik}
  &=\operatorname{Cov}(X_I(t),X_k(t))                                      \\
  &=\int_0^\infty e^{r\mathcal{A}_{II}}a_{Ik}
     \mathbb{E}\bigl[X_k(t-r)X_k(t)\bigr] \,dr                              \\
  &=\Sigma_{\lambda,kk}
    \int_0^\infty e^{-\lambda r}e^{r\mathcal{A}_{II}}a_{Ik}\,dr              \\
  &=\frac{\varepsilon_0}{2\lambda^2}
    (\lambda I-\mathcal{A}_{II})^{-1}a_{Ik}.
\end{align*}
Equivalently, the $Ik$ block of the Lyapunov equation is
\[
  (\mathcal{A}_{II}-\lambda I)\Sigma_{\lambda,Ik}
  +a_{Ik}\Sigma_{\lambda,kk}=0,
\]
which rearranges to the same positive-sign resolvent formula.
\end{proof}

\subsection{Cameron--Martin energy}

Let $Q_0\coloneqq\Sigma_{II}^{(0)}$ be the background bulk covariance, and let
$H_{Q_0}\coloneqq\Range(Q_0^{1/2})$ have the energy norm
$|u|_{Q_0}=\|Q_0^{-1/2}u\|_{\mathcal{H}_I}$.  Assumption~\ref{assum:cm_stability}
states that $a_{Ik}\in H_{Q_0}$ and that the bulk resolvent is bounded on this
energy space:
\begin{equation}
  |(\lambda I-\mathcal{A}_{II})^{-1}u|_{Q_0}
  \le \frac{M}{\lambda+\omega}|u|_{Q_0}.
  \label{eq:cm_resolvent}
\end{equation}

\begin{lemma}[Finite-energy leakage]
\label{lem:finite_energy_leakage}
The cross covariance belongs to $H_{Q_0}$ and satisfies
\begin{equation}
 |\Sigma_{\lambda,Ik}|_{Q_0}
 \le
 \frac{\varepsilon_0M}{2\lambda^2(\lambda+\omega)}
 |a_{Ik}|_{Q_0}
 =O(\lambda^{-3}).
 \label{eq:cross_energy}
\end{equation}
\end{lemma}

\begin{proof}
Apply \eqref{eq:cm_resolvent} to the resolvent identity
\eqref{eq:SigmaIk_exact_vector}.  This proves both the range statement and the
energy bound.
\end{proof}

\begin{lemma}[Form-domain energy domination]
\label{lem:form_domain_energy_domination}
Let $C$ and $Q$ be bounded non-negative self-adjoint operators on a Hilbert
space and assume $C\succeq Q$.  Then
\begin{equation}
  H_Q=\Range(Q^{1/2})\subseteq \Range(C^{1/2})=H_C,
  \label{eq:range_inclusion_CQ}
\end{equation}
and for every $v\in H_Q$,
\begin{equation}
  |v|_C\le |v|_Q.
  \label{eq:energy_domination_CQ}
\end{equation}
In particular, with $C=\Sigma_{\lambda,II}$ and $Q=Q_0$,
\begin{equation}
  |\Sigma_{\lambda,II}^{-1/2}v|
  \le |Q_0^{-1/2}v|,
  \qquad v\in H_{Q_0}.
  \label{eq:Sigma_energy_dominated_by_Q0}
\end{equation}
\end{lemma}

\begin{proof}
The Loewner inequality $Q\le C$ is equivalent to
\[
  Q^{1/2}(Q^{1/2})^*\le C^{1/2}(C^{1/2})^*.
\]
Douglas' range inclusion theorem
therefore gives a contraction $K$ on the ambient Hilbert space such that
\begin{equation}
  Q^{1/2}=C^{1/2}K,
  \qquad \|K\|\le1.
  \label{eq:douglas_factorization}
\end{equation}
This proves \eqref{eq:range_inclusion_CQ}.  If $v\in H_Q$, write
$v=Q^{1/2}x$ with $x$ chosen as the minimal-norm preimage, so
$\|x\|=|v|_Q$.  By \eqref{eq:douglas_factorization},
$v=C^{1/2}Kx$.  The Cameron--Martin norm $|v|_C$ is the minimal norm among all
$y$ with $C^{1/2}y=v$.  Hence
\[
  |v|_C\le \|Kx\|\le\|x\|=|v|_Q.
\]
Taking $C=\Sigma_{\lambda,II}$ and $Q=Q_0$ is justified by
Lemma~\ref{lem:bulk_domination_convolution}, and gives
\eqref{eq:Sigma_energy_dominated_by_Q0}.
\end{proof}

\section{Anderson--Trapp Shorting and the Rank-One Residual}
\label{app:shorting_residual}

In finite dimensions, the Schur complement is only the bounded-inverse
shadow of Anderson--Trapp shorting.  The Hilbert space argument therefore
starts from shorting and recovers the Schur expression only when the
complementary block is boundedly invertible.

\begin{lemma}[Schur complement as the bounded-inverse shadow of shorting]
\label{lem:schur_shadow_shorting}
Let $\mathcal{E}$ and $\mathcal{F}$ be Hilbert spaces and let
\[
  T=\begin{pmatrix}A&B^*\\B&D\end{pmatrix}\ge0
\]
act on $\mathcal{E}\oplus\mathcal{F}$.  If $D$ is boundedly invertible on
$\mathcal{F}$, then the Anderson--Trapp short of $T$ to $\mathcal{E}$ has the
quadratic form
\begin{equation}
  \langle x,\mathcal{S}_{\mathcal{E}}(T)x\rangle
  =
  \langle x,Ax\rangle-\langle Bx,D^{-1}Bx\rangle.
  \label{eq:short_shadow}
\end{equation}
Thus $B^*D^{-1}B$ is exactly the variational shorting subtraction in the special
case where $D^{-1}$ is a bounded operator.
\end{lemma}

\begin{proof}
By the variational definition of the shorted operator,
\[
  \langle x,\mathcal{S}_{\mathcal{E}}(T)x\rangle
  =
  \inf_{y\in\mathcal{F}}
  \left\langle
  \begin{pmatrix}x\\y\end{pmatrix},
  T
  \begin{pmatrix}x\\y\end{pmatrix}
  \right\rangle .
\]
The displayed quadratic form equals
$\langle x,Ax\rangle+2\operatorname{Re}\langle Bx,y\rangle+
\langle y,Dy\rangle$.  Since $D$ is boundedly invertible and positive, the unique
minimizer is $y=-D^{-1}Bx$.  Substitution gives \eqref{eq:short_shadow}.
\end{proof}

\begin{lemma}[Rayleigh quotient representation of covariance energy]
\label{lem:rayleigh_covariance_energy}
Let $C$ be a bounded non-negative self-adjoint operator on a Hilbert space
$\mathcal{F}$ and let $z\in\mathcal{F}$.  Then
\begin{equation}
  \sup_{\langle u,Cu\rangle>0}
  \frac{|\langle u,z\rangle|^2}{\langle u,Cu\rangle}
  =
  \begin{cases}
    |z|_C^2, & z\in H_C,\\
    +\infty, & z\notin H_C,
  \end{cases}
  \label{eq:rayleigh_equals_energy}
\end{equation}
where $H_C=\Range(C^{1/2})$ and $|z|_C=\|C^{-1/2}z\|$ is the
Cameron--Martin energy norm, interpreted as the minimal preimage norm.
\end{lemma}

\begin{proof}
By the spectral theorem, write $C e_j=\gamma_j e_j$ on the positive spectral
subspace and expand $z=\sum_j z_j e_j+z_0$ with $z_0\in\ker C$.  If $z_0\ne0$,
then positivity of the denominator is unaffected by adding large multiples of
$z_0$ to $u$, and the supremum is infinite.  Thus finiteness forces
$z\perp\ker C$.  On the positive spectral subspace,
\[
  \langle u,Cu\rangle=\sum_j\gamma_j|u_j|^2,
  \qquad
  \langle u,z\rangle=\sum_j\overline{u_j}z_j.
\]
Cauchy's inequality gives
\[
  |\langle u,z\rangle|^2
  \le
  \left(\sum_j\gamma_j|u_j|^2\right)
  \left(\sum_j\frac{|z_j|^2}{\gamma_j}\right).
\]
If $\sum_j|z_j|^2/\gamma_j=\infty$, truncating to the first $m$ positive modes
shows that the supremum diverges.  If this sum is finite, choose
$u_j^{(m)}=z_j/\gamma_j$ on the first $m$ positive modes and zero otherwise;
the corresponding Rayleigh quotients increase to
$\sum_j|z_j|^2/\gamma_j=|z|_C^2$.  This proves
\eqref{eq:rayleigh_equals_energy}.
\end{proof}

For the covariance block $\Sigma_\lambda$ on
$\mathbb{C}\phi_k\oplus\mathcal{H}_I$, the shorting subtraction is the
variational quantity
\begin{equation}
 R_\lambda
 =
 \sup_{\langle u,\Sigma_{\lambda,II}u\rangle>0}
 \frac{|\langle u,\Sigma_{\lambda,Ik}\rangle|^2}
      {\langle u,\Sigma_{\lambda,II}u\rangle}.
 \label{eq:schur_variational}
\end{equation}
The null directions of the denominator do not change the supremum for a
non-negative block covariance: positivity of the full block forces the
cross-vector to annihilate those null directions.

\begin{proposition}[Rank-one shorted residual]
\label{prop:rank_one_shorted_residual}
Under Assumption~\ref{assum:core_analytical},
\begin{equation}
  S_{k,\lambda}
  =\Sigma_{\lambda,kk}-R_\lambda,
  \qquad
  R_\lambda
  =|\Sigma_{\lambda,Ik}|_{\Sigma_{\lambda,II}}^2
  =\|\Sigma_{\lambda,Ik}\|_{H_{\Sigma_{\lambda,II}}}^2,
  \label{eq:shorting_residual_energy_identity}
\end{equation}
and
\begin{equation}
  0\le R_\lambda
  \le |\Sigma_{\lambda,Ik}|_{Q_0}^2
  =O(\lambda^{-6}).
  \label{eq:shorted_residual_bound}
\end{equation}
Consequently
\begin{equation}
  S_{k,\lambda}
  =
  \frac{\varepsilon_0}{2\lambda^2}-O(\lambda^{-6}),
  \qquad
  m_2^{\mathrm{Schur}}=\frac{\varepsilon_0}{2},
  \label{eq:Schur_asympt}
\end{equation}
and
\[
  S_{k,\lambda}^{-1}
  =
  \frac{2}{\varepsilon_0}\lambda^2+O(\lambda^{-2}).
\]
\end{proposition}

\begin{proof}
The first identity $S_{k,\lambda}=\Sigma_{\lambda,kk}-R_\lambda$ is the
one-dimensional Anderson--Trapp short of the positive block covariance
$\Sigma_\lambda$ to the target coordinate.  Applying
Lemma~\ref{lem:rayleigh_covariance_energy} with
$C=\Sigma_{\lambda,II}$ and $z=\Sigma_{\lambda,Ik}$ gives the exact identity
\[
  R_\lambda=|\Sigma_{\lambda,Ik}|_{\Sigma_{\lambda,II}}^2,
\]
provided the vector lies in $H_{\Sigma_{\lambda,II}}$.

This range statement follows from the previous section.  Specifically, Lemma~\ref{lem:bulk_domination_convolution} gives
$\Sigma_{\lambda,II}\succeq Q_0$, while Lemma~\ref{lem:form_domain_energy_domination} shows that
$H_{Q_0}\subseteq H_{\Sigma_{\lambda,II}}$ with
\[
  |v|_{\Sigma_{\lambda,II}}\le |v|_{Q_0},
  \qquad v\in H_{Q_0}.
\]
Since Lemma~\ref{lem:finite_energy_leakage} gives
$\Sigma_{\lambda,Ik}\in H_{Q_0}$ and
$|\Sigma_{\lambda,Ik}|_{Q_0}=O(\lambda^{-3})$, we obtain
\[
  0\le R_\lambda
  =|\Sigma_{\lambda,Ik}|_{\Sigma_{\lambda,II}}^2
  \le |\Sigma_{\lambda,Ik}|_{Q_0}^2
  =O(\lambda^{-6}).
\]
Combining this with \eqref{eq:Sigmakk_exact} proves
\eqref{eq:Schur_asympt}.  Finally,
\[
  S_{k,\lambda}
  =\frac{\varepsilon_0}{2\lambda^2}
   \left(1-\frac{2\lambda^2}{\varepsilon_0}R_\lambda\right),
  \qquad
  \frac{2\lambda^2}{\varepsilon_0}R_\lambda=O(\lambda^{-4}),
\]
so the reciprocal expansion follows from the scalar Neumann expansion.
\end{proof}

\section{Galerkin Approximation and the Failure of Direct Inversion}
\label{app:galerkin_shorting}

Galerkin approximation is used here only to approximate trace-class covariances
and nested variational shorting problems.  It is not a license to replace the
Hilbert-space short by a global inverse of a compact covariance block.

\begin{lemma}[Trace-norm convergence of Galerkin covariances]
\label{lem:galerkin_trace_norm}
Let $\Pi_n:\mathcal{H}\to\mathcal{H}_n$ be finite-rank orthogonal projections
with $\Pi_n\phi_k=\phi_k$, $\Pi_n\mathcal{H}_I\subset\mathcal{H}_I$, and
$\Pi_n\to I$ strongly.  For fixed $\lambda>0$, set
\[
  \mathcal{A}_n^{(\lambda)}
  =\Pi_n\mathcal{A}^{(\lambda)}|_{\mathcal{H}_n},
  \qquad
  \mathcal{D}_n^{(\lambda)}
  =\Pi_n\mathcal{D}^{(\lambda)}\Pi_n .
\]
Assume that $e^{t\mathcal{A}_n^{(\lambda)}}x\to
e^{t\mathcal{A}^{(\lambda)}}x$ for every $x$, uniformly for $t$ in compact
intervals, that the semigroups are uniformly exponentially stable, satisfying
\[
  \|e^{t\mathcal{A}_n^{(\lambda)}}\|,\
  \|e^{t\mathcal{A}^{(\lambda)}}\|\le Me^{-\beta t},
\]
and that $\mathcal{D}_n^{(\lambda)}\to\mathcal{D}^{(\lambda)}$ in
$\mathcal{S}_1(\mathcal{H})$.  Then
\[
  \Sigma_{n,\lambda}
  =
  \int_0^\infty
  e^{t\mathcal{A}_n^{(\lambda)}}
  \mathcal{D}_n^{(\lambda)}
  e^{t\mathcal{A}_n^{(\lambda)*}}\,dt
  \longrightarrow
  \Sigma_\lambda
  \quad\text{in }\mathcal{S}_1(\mathcal{H}).
\]
\end{lemma}

\begin{proof}
Omitting the superscript $(\lambda)$, write the difference as a Bochner
integral and split the integrand:
\begin{multline*}
  \|e^{t\mathcal{A}_n}\mathcal{D}_ne^{t\mathcal{A}_n^*}
    -e^{t\mathcal{A}}\mathcal{D}e^{t\mathcal{A}^*}\|_{\mathcal{S}_1}
  \le
  M^2e^{-2\beta t}\|\mathcal{D}_n-\mathcal{D}\|_{\mathcal{S}_1}\\
  +
  \|e^{t\mathcal{A}_n}\mathcal{D}e^{t\mathcal{A}_n^*}
    -e^{t\mathcal{A}}\mathcal{D}e^{t\mathcal{A}^*}\|_{\mathcal{S}_1}.
\end{multline*}
The first term tends to zero.  For the second, Gr\"umm's theorem for trace
ideals implies trace-norm convergence from strong convergence of the left and
right semigroups, uniform boundedness, and $\mathcal{D}\in\mathcal{S}_1$
\cite{simon2005trace,daprato2014stochastic}.  The full integrand is dominated by
\[
  M^2e^{-2\beta t}
  \bigl(\|\mathcal{D}_n\|_{\mathcal{S}_1}
        +\|\mathcal{D}\|_{\mathcal{S}_1}\bigr),
\]
with a finite uniform trace-norm bound because
$\mathcal{D}_n\to\mathcal{D}$ in $\mathcal{S}_1$.  Dominated convergence for
Bochner integrals gives the claim.
\end{proof}

\begin{lemma}[Spectral truncation energy bound]
\label{lem:galerkin_finite_energy_monotonicity}
Assume that the Galerkin projections are spectral truncations of the background
bulk covariance.  Thus, on $\mathcal{H}_I$,
\[
  Q_0e_j=q_je_j,
  \qquad q_j>0,
  \qquad
  \Pi_n^I u=\sum_{j=1}^n \langle u,e_j\rangle e_j.
\]
Let $Q_0^{(n)}=\Pi_n^I Q_0\Pi_n^I|_{\Pi_n^I\mathcal{H}_I}$ and
$a_{Ik}^{(n)}=\Pi_n^I a_{Ik}$.  If $a_{Ik}\in H_{Q_0}$, then
\begin{equation}
  |a_{Ik}^{(n)}|_{Q_0^{(n)}}^2
  =
  \sum_{j=1}^n\frac{|\langle a_{Ik},e_j\rangle|^2}{q_j}
  \uparrow
  |a_{Ik}|_{Q_0}^2.
  \label{eq:galerkin_energy_monotone}
\end{equation}
In particular,
\begin{equation}
  \sup_n |a_{Ik}^{(n)}|_{Q_0^{(n)}}
  \le |a_{Ik}|_{Q_0}<\infty.
  \label{eq:galerkin_uniform_finite_energy}
\end{equation}
\end{lemma}

\begin{proof}
Writing $a_j=\langle a_{Ik},e_j\rangle$, the Cameron--Martin norm of the
full vector is
\[
  |a_{Ik}|_{Q_0}^2=\sum_{j=1}^\infty\frac{|a_j|^2}{q_j}<\infty.
\]
The truncated covariance has eigenvalues $q_1,\ldots,q_n$ on
$\Pi_n^I\mathcal{H}_I$, and $a_{Ik}^{(n)}=\sum_{j=1}^n a_je_j$.  Therefore
\eqref{eq:galerkin_energy_monotone} follows directly from the spectral formula
for the finite-dimensional Cameron--Martin norm.  The partial sums are monotone
and bounded by the full sum, giving \eqref{eq:galerkin_uniform_finite_energy}.
\end{proof}

\begin{lemma}[Galerkin covariance residuals and Rayleigh--Ritz shorting]
\label{lem:galerkin_variational_shorting}
For every block-preserving Galerkin truncation satisfying the hypotheses of
Lemma~\ref{lem:galerkin_trace_norm} and the uniform finite-energy condition
\[
  a_{Ik}^{(n)}\in H_{Q_0^{(n)}},
  \qquad
  \sup_n |a_{Ik}^{(n)}|_{Q_0^{(n)}}<\infty,
\]
the truncated residual satisfies
\[
  S_{k,\lambda}^{(n)}
  =
  \frac{\varepsilon_0}{2\lambda^2}-R_{n,\lambda},
  \qquad
  0\le R_{n,\lambda}\le C\lambda^{-6},
\]
with $C$ independent of $n$.  Hence
$m_2^{\mathrm{Schur},(n)}=\varepsilon_0/2$ for every admissible truncation.
No monotonicity claim is made for the family $R_{n,\lambda}$.

Separately, for fixed $\lambda$, nested finite-dimensional test spaces in the
bulk Cameron--Martin space define Rayleigh--Ritz restrictions
$R_\lambda^{[n]}$ of the fixed infinite-dimensional variational shorting
problem.  These satisfy
\[
  R_\lambda^{[n]}\uparrow R_\lambda,
  \qquad
  S_{k,\lambda}^{[n]}\coloneqq\Sigma_{\lambda,kk}-R_\lambda^{[n]}
  \downarrow S_{k,\lambda}.
\]
\end{lemma}

\begin{proof}
The truncated operators have the same triangular form as
\eqref{eq:exact_blocks}, so the finite-dimensional Lyapunov blocks yield
\[
  \Sigma_{n,\lambda,kk}=\frac{\varepsilon_0}{2\lambda^2},
  \qquad
  \Sigma_{n,\lambda,Ik}
  =
  \frac{\varepsilon_0}{2\lambda^2}
  (\lambda I-\mathcal{A}_{II}^{(n)})^{-1}a_{Ik}^{(n)}.
\]
The uniform finite-energy assumption and the uniform $H_{Q_0^{(n)}}$ resolvent
bound give
$|\Sigma_{n,\lambda,Ik}|_{Q_0^{(n)}}=O(\lambda^{-3})$ uniformly in $n$.
The same Douglas domination argument used in
Proposition~\ref{prop:rank_one_shorted_residual} then yields
$0\le R_{n,\lambda}\le C\lambda^{-6}$ and the stated value of
$m_2^{\mathrm{Schur},(n)}$.  For spectral truncations
$a_{Ik}^{(n)}=\Pi_n^I a_{Ik}$, the required uniform finite-energy condition is
provided by Lemma~\ref{lem:galerkin_finite_energy_monotonicity}.

For the fixed-problem Rayleigh--Ritz statement, let
$C_\lambda=\Sigma_{\lambda,II}$ and let $V_m$ be nested finite-dimensional
subspaces whose union is dense in $H_{C_\lambda}=\Range(C_\lambda^{1/2})$.
Then
\[
  R_\lambda^{[m]}
  =
  \sup_{\substack{u\in V_m\\ \langle u,C_\lambda u\rangle>0}}
  \frac{|\langle u,\Sigma_{\lambda,Ik}\rangle|^2}
       {\langle u,C_\lambda u\rangle}
\]
is monotone increasing in $m$ and bounded above by $R_\lambda$.  Density in the
Cameron--Martin energy norm gives $R_\lambda^{[m]}\uparrow R_\lambda$ and
$S_{k,\lambda}^{[m]}=\Sigma_{\lambda,kk}-R_\lambda^{[m]}\downarrow
S_{k,\lambda}$.  In an eigenbasis of $C_\lambda$, this is the Rayleigh--Ritz
identity
\[
  R_\lambda^{[m]}
  =
  \sum_{j=1}^m
  \frac{|\langle\Sigma_{\lambda,Ik},e_j\rangle|^2}{\sigma_j}
  \uparrow
  \|\Sigma_{\lambda,Ik}\|_{H_{C_\lambda}}^2.
\]
\end{proof}


The obstruction to direct inversion is not numerical bookkeeping.  In the GFF
spectral regime it can be stated as a precise spectral theorem.

\begin{lemma}[GFF spectral obstruction to direct inversion]
\label{lem:gff_spectral_obstruction_appendix}
Assume that the background covariance and the bulk drift are simultaneously
represented in a bulk eigenbasis $(e_j)_{j\ge1}$ such that
\[
  Q_0e_j=\sigma_je_j,
  \qquad
  \sigma_j\asymp j^{-2},
\]
and the positive damping eigenvalues of $-\mathcal{A}_{II}$ satisfy
\[
  \mu_j\asymp j^2 .
\]
Let $Q_0^{(N)}$ be the $N$-mode Galerkin truncation in this basis.  Then:
\begin{enumerate}[label=(\roman*)]
\item $Q_0$ is compact and injective with $0$ in its spectrum.  Hence
$Q_0^{-1}$ is not a bounded ambient Hilbert-space operator, and a formula
containing the global inverse of the bulk covariance block is not an operator
identity on the ambient space.
\item
\[
  \kappa(Q_0^{(N)})
  =
  \frac{\max_{1\le j\le N}\sigma_j}{\min_{1\le j\le N}\sigma_j}
  \asymp N^2 .
\]
Thus direct Galerkin inversion amplifies errors at the inverse GFF precision
scale.
\item The diagonal resolvent filter $(\lambda+\mu_j)^{-1}$ has a moving critical
mode $j^*(\lambda)$ determined by $\mu_{j^*(\lambda)}\asymp\lambda$, and hence
\[
  j^*(\lambda)\asymp\sqrt{\lambda}.
\]
For $j\ll j^*(\lambda)$ one has
$(\lambda+\mu_j)^{-1}\asymp\lambda^{-1}$, whereas for
$j\gg j^*(\lambda)$ one has
$(\lambda+\mu_j)^{-1}\asymp\mu_j^{-1}\asymp j^{-2}$.  Therefore a fixed
Galerkin dimension $N$ does not resolve the large-$\lambda$ transition unless
$N$ grows on the moving scale $\sqrt{\lambda}$.
\end{enumerate}
Consequently, the infinite-dimensional object that survives the singular limit
is the Anderson--Trapp variational short, not the direct inverse of a compact
bulk covariance block.
\end{lemma}

\begin{proof}
The first claim follows from the spectral theorem: $\sigma_j>0$ and
$\sigma_j\to0$, so $Q_0$ is injective and compact, while $0$ is a spectral point.
The inverse is defined only on its range as an unbounded operator, equivalently
as the Cameron--Martin energy form.

For the condition number, the estimate $\sigma_j\asymp j^{-2}$ gives
$\max_{1\le j\le N}\sigma_j\asymp\sigma_1\asymp1$. Meanwhile, we have
$\min_{1\le j\le N}\sigma_j\asymp\sigma_N\asymp N^{-2}$, proving that
$\kappa(Q_0^{(N)})\asymp N^2$.

For the critical mode, the transition of $(\lambda+\mu_j)^{-1}$ occurs when the
two summands in the denominator are comparable.  Since $\mu_j\asymp j^2$, the
relation $\mu_j\asymp\lambda$ is equivalent to $j\asymp\sqrt{\lambda}$.  The two
resolvent regimes follow by comparing $\mu_j$ with $\lambda$ below and above this
scale.
\end{proof}

\section{Gaussian Radon Conditioning and the Fredholm Finite Part}
\label{app:gaussian_fredholm}

The Gaussian and Fredholm layer is the Radon representation of the shorted
residual constructed above.  It is not a second mechanism: disintegration
supplies the exact-evidence law, while the shorted covariance residual supplies
the finite conditional variance used by the Gaussian entropy calculation.

\subsection{Pointwise disintegration}

Let $\mu_\lambda=\mathcal{N}(m_\lambda,\Sigma_\lambda)$ be a Radon Gaussian
measure on $\mathcal{H}=\mathcal{H}_k\oplus\mathcal{H}_I$, and let
$\pi_k(X)=\langle X,\phi_k\rangle$.  Write
\[
  \sigma_\lambda=\Sigma_{\lambda,kk},
  \qquad
  C_\lambda=\Sigma_{\lambda,II},
  \qquad
  z_\lambda=\Sigma_{\lambda,Ik}.
\]

\begin{lemma}[Gaussian disintegration covariance]
\label{lem:disintegration}
For each $t\in\mathbb{R}$, there is a unique weakly continuous Gaussian
disintegration member
$\mu_\lambda^t=\mathcal{N}(m_\lambda^t,\Sigma_\lambda^t)$ supported on
$X_k=t$, where
\[
  m_\lambda^t=
  \begin{pmatrix}
    t\\
    m_{\lambda,I}
    +z_\lambda\sigma_\lambda^{-1}(t-m_{\lambda,k})
  \end{pmatrix},
\]
and
\begin{equation}
  \Sigma_\lambda^t=
  \begin{pmatrix}
    0&0\\
    0&C_\lambda-z_\lambda\sigma_\lambda^{-1}z_\lambda^*
  \end{pmatrix}.
  \label{eq:conditional_covariance_block}
\end{equation}
The bulk conditional covariance
\begin{equation}
  C_\lambda^t
  \coloneqq
  C_\lambda-z_\lambda\sigma_\lambda^{-1}z_\lambda^*
  \label{eq:bulk_conditional_covariance}
\end{equation}
is a positive trace-class covariance operator on $\mathcal{H}_I$.
At $t=x_k^*$, the observed law $\mu_{\mathrm{obs}}=\mu_\lambda^{x_k^*}$
satisfies
\[
  \mathbb{E}_{\mathrm{obs}}[X_k]-x_k^*=0,
  \qquad
  \operatorname{Var}_{\mathrm{obs}}(X_k)=0.
\]
\end{lemma}

\begin{proof}
The map $\pi_k$ is a continuous linear map to $\mathbb{R}$, and its marginal
variance is $\sigma_\lambda>0$.  Regular conditional probabilities therefore
exist because $\mathcal{H}$ is Polish.  To identify the continuous Gaussian
version, center the variables and set
\[
  \widetilde X_I
  =X_I-m_{\lambda,I}-z_\lambda\sigma_\lambda^{-1}(X_k-m_{\lambda,k}).
\]
For every $h\in\mathcal{H}_I$,
\[
  \operatorname{Cov}\bigl(\langle h,\widetilde X_I\rangle,
  X_k-m_{\lambda,k}\bigr)
  =\langle h,z_\lambda\rangle
   -\langle h,z_\lambda\rangle\sigma_\lambda^{-1}\sigma_\lambda
  =0.
\]
The pair is jointly Gaussian, hence $\widetilde X_I$ is independent of $X_k$.
Conditioning on $X_k=t$ therefore shifts only the deterministic affine part and
leaves the covariance of $\widetilde X_I$ unchanged.  Direct covariance
calculation gives
\[
  \operatorname{Cov}(\widetilde X_I)
  =C_\lambda-z_\lambda\sigma_\lambda^{-1}z_\lambda^*.
\]
This proves the displayed conditional mean and covariance.  The operator
$C_\lambda^t$ is positive because it is a covariance.  It is trace-class because
$C_\lambda$ is trace-class and $z_\lambda\sigma_\lambda^{-1}z_\lambda^*$ is
rank one.  The affine formulas are continuous in $t$, so this gives the unique
weakly continuous disintegration member along the scalar observation map
\cite{bogachev1998gaussian,lagatta2013continuous,steinwart2024conditioning,chen2024gaussian}.  Evaluating at $x_k^*$ fixes
the target coordinate almost surely.
\end{proof}

\subsection{Gaussian energy projection and entropy}

Since $C_\lambda$ is trace-class, its inverse is not a bounded ambient operator;
the relevant object is the Cameron--Martin space
$H_{C_\lambda}=\Range(C_\lambda^{1/2})$.

\begin{lemma}[Measurable Cameron--Martin energy projection]
\label{lem:conditional_predictor}
The cross covariance belongs to $H_{C_\lambda}$ and
\[
  \|z_\lambda\|_{H_{C_\lambda}}
  \le |z_\lambda|_{Q_0}=O(\lambda^{-3}).
\]
The conditional mean projection is the measurable Wiener functional
\[
  \mu_{k|I}(X_I)
  =
  m_{\lambda,k}
  +W_{z_\lambda}(X_I-m_{\lambda,I}),
\]
and, for all sufficiently large $\lambda$,
\[
  \operatorname{Var}(X_k\mid X_I)
  =
  \sigma_\lambda-\|z_\lambda\|_{H_{C_\lambda}}^2
  =S_{k,\lambda}>0.
\]
\end{lemma}

\begin{proof}
The domination $C_\lambda=\Sigma_{\lambda,II}\succeq Q_0$ and
Lemma~\ref{lem:form_domain_energy_domination} give
$H_{Q_0}\subseteq H_{C_\lambda}$ with contraction of energy norms.  Combining
this inclusion with \eqref{eq:cross_energy} gives the norm bound.  The range
condition $z_\lambda\in H_{C_\lambda}$ is exactly the condition for the
measurable Wiener functional associated with $z_\lambda$ to exist under the bulk
Gaussian law \cite{bogachev1998gaussian,steinwart2024conditioning}.  The variance identity is the
one-dimensional shorting identity from Proposition~\ref{prop:rank_one_shorted_residual};
positivity holds for large $\lambda$ because
$S_{k,\lambda}=\varepsilon_0/(2\lambda^2)-O(\lambda^{-6})$.
\end{proof}

\begin{lemma}[Feldman--Hajek rank-one perturbation]
\label{lem:feldman_hajek}
There is $\lambda_0>0$ such that for $\lambda\ge\lambda_0$ the bulk marginals
$\mu_{\mathrm{obs},I}$ and $\mu_{\lambda,I}$ are equivalent.  The relative
covariance perturbation has at most one nontrivial eigenvalue, namely
\begin{equation}
  1-\rho_\lambda,
  \qquad
  \rho_\lambda
  =
  \frac{\|z_\lambda\|_{H_{C_\lambda}}^2}{\sigma_\lambda}
  =
  1-\frac{S_{k,\lambda}}{\sigma_\lambda}
  =O(\lambda^{-4}).
  \label{eq:rho_lambda_rank_one}
\end{equation}
Moreover, if $\Delta_\lambda=x_k^*-m_{\lambda,k}$, then
\begin{equation}
  2\KL(\mu_{\mathrm{obs},I}\parallel\mu_{\lambda,I})
  =
  \frac{\Delta_\lambda^2}{\sigma_\lambda}\rho_\lambda
  -\log(1-\rho_\lambda)-\rho_\lambda.
  \label{eq:bulk_KL_rank_one}
\end{equation}
In particular,
$\KL(\mu_{\mathrm{obs},I}\parallel\mu_{\lambda,I})=O(\lambda^{-4})$.
\end{lemma}

\begin{proof}
Put $g_\lambda=C_\lambda^{-1/2}z_\lambda$ and
$\alpha_\lambda=\|g_\lambda\|^2=\|z_\lambda\|_{H_{C_\lambda}}^2$.  If
$z_\lambda=0$, then $\rho_\lambda=0$ and the covariance statement is trivial.
Otherwise let $e_\lambda=g_\lambda/\|g_\lambda\|$.  The conditional bulk
covariance from \eqref{eq:bulk_conditional_covariance} can be written in
relative form as
\begin{equation}
  C_\lambda^{x_k^*}
  =C_\lambda^{1/2}
    (I-\rho_\lambda e_\lambda\otimes e_\lambda)
    C_\lambda^{1/2},
  \qquad
  \rho_\lambda=\alpha_\lambda/\sigma_\lambda.
  \label{eq:relative_rank_one_covariance}
\end{equation}
Thus the relative covariance is the identity on
$e_\lambda^\perp$ and has the single nontrivial eigenvalue
$1-\rho_\lambda$ along $e_\lambda$.

By Proposition~\ref{prop:rank_one_shorted_residual},
$\alpha_\lambda=R_\lambda=O(\lambda^{-6})$, while
$\sigma_\lambda=\varepsilon_0/(2\lambda^2)$.  Hence
$\rho_\lambda=O(\lambda^{-4})$ and, for all sufficiently large $\lambda$,
$0\le\rho_\lambda<1$.  Then
$I-\rho_\lambda e_\lambda\otimes e_\lambda$ is boundedly positive and invertible,
so the Cameron--Martin ranges of $C_\lambda^{x_k^*}$ and $C_\lambda$ coincide.
The covariance perturbation is rank one, hence Hilbert--Schmidt in relative
coordinates.  The mean shift is
$z_\lambda\sigma_\lambda^{-1}\Delta_\lambda\in H_{C_\lambda}$ and has
$C_\lambda$-energy squared
\[
  \bigl\|C_\lambda^{-1/2}z_\lambda\sigma_\lambda^{-1}\Delta_\lambda\bigr\|^2
  =\frac{\Delta_\lambda^2}{\sigma_\lambda}\rho_\lambda.
\]
Feldman--Hajek therefore gives equivalence of the two bulk Gaussian laws; recent Hilbert-space Gaussian inverse-problem work uses the same measure-class condition for finite-rank posterior updates \cite{carerelie2024covariance,carerelie2025mean}.

The Gaussian relative entropy formula for equivalent Gaussian measures gives
one contribution from the mean and one from the nontrivial relative covariance
eigenvalue:
\[
  2\KL
  =
  \frac{\Delta_\lambda^2}{\sigma_\lambda}\rho_\lambda
  +\bigl[(1-\rho_\lambda)-1-\log(1-\rho_\lambda)\bigr],
\]
which is \eqref{eq:bulk_KL_rank_one}; see also recent Hilbert-space log-determinant/KL treatments \cite{quang2022kl,carerelie2024covariance}.  Since
$\Delta_\lambda=O(\lambda^{-1})$, $\sigma_\lambda=O(\lambda^{-2})$, and
$\rho_\lambda=O(\lambda^{-4})$, the mean contribution is $O(\lambda^{-4})$;
the logarithmic remainder is $O(\rho_\lambda^2)$ after cancellation of the
linear term.  Thus the bulk entropy is $O(\lambda^{-4})$.
\end{proof}

\subsection{Mollified divergence and finite part}

For $\delta>0$, set
\[
  \nu_{\lambda,\delta}
  =
  \mu_{\mathrm{obs},I}\otimes\mathcal{N}(x_k^*,\delta^2)
\]
under the product identification
$\mathcal{H}=\mathcal{H}_I\times\mathcal{H}_k$.

\begin{proposition}[Mollification limit]
\label{prop:mollification_limit}
For $\lambda\ge\lambda_0$ as in Lemma~\ref{lem:feldman_hajek}, define
\[
  C_{\lambda,\delta}
  =
  \frac12\left(\log\frac{S_{k,\lambda}}{\delta^2}-1\right).
\]
Then
\[
  \lim_{\delta\downarrow0}
  \left[
    \KL(\nu_{\lambda,\delta}\parallel\mu_\lambda)
    -C_{\lambda,\delta}
  \right]
  =
  \mathfrak{J}_\lambda(\mu_{\mathrm{obs}};\mu_\lambda).
\]
\end{proposition}

\begin{proof}
Let $r_\lambda(X)=X_k-\mu_{k|I}(X_I)$.  The entropy chain rule and the scalar
Gaussian formula give
\[
  \KL(\nu_{\lambda,\delta}\parallel\mu_\lambda)
  =
  \KL(\mu_{\mathrm{obs},I}\parallel\mu_{\lambda,I})
  +C_{\lambda,\delta}
  +\frac{\delta^2}{2S_{k,\lambda}}
  +\frac{1}{2S_{k,\lambda}}
   \mathbb{E}_{\mu_{\mathrm{obs}}}[r_\lambda(X)^2].
\]
The $\delta^2$ term vanishes as $\delta\downarrow0$, and the remaining terms
are exactly Definition~\ref{def:unified_action}.
\end{proof}

\begin{lemma}[KL finite-part cancellation algebra]
\label{lem:kl_finite_part_algebra}
Let $t\in\mathbb{R}$, let $\mu_\lambda^t$ be the Gaussian disintegration member,
and set $\Delta_t=t-m_{\lambda,k}$.  For $\lambda\ge\lambda_0$,
\begin{equation}
  \KL(\mu_{\lambda,I}^t\parallel\mu_{\lambda,I})
  =
  \frac12\frac{\Delta_t^2}{\sigma_\lambda}\rho_\lambda
  -\frac12\log(1-\rho_\lambda)-\frac12\rho_\lambda,
  \label{eq:bulk_KL_t_algebra}
\end{equation}
and
\begin{equation}
  \frac{1}{2S_{k,\lambda}}
  \mathbb{E}_{\mu_\lambda^t}
  \left[(X_k-\mu_{k|I}(X_I))^2\right]
  =
  \frac12\rho_\lambda
  +\frac12\frac{\Delta_t^2}{\sigma_\lambda}(1-\rho_\lambda).
  \label{eq:residual_t_algebra}
\end{equation}
Consequently,
\begin{equation}
  \mathfrak{J}_\lambda(t)
  =
  \frac{\Delta_t^2}{2\sigma_\lambda}
  -\frac12\log(1-\rho_\lambda).
  \label{eq:finite_part_exact_rank_one}
\end{equation}
\end{lemma}

\begin{proof}
Equation \eqref{eq:bulk_KL_t_algebra} is
\eqref{eq:bulk_KL_rank_one} with $x_k^*$ replaced by $t$; this is the rank-one specialization of the standard equivalent-Gaussian KL formula \cite{bogachev1998gaussian,quang2022kl}.
It remains to compute the second term.  Under $\mu_\lambda^t$, the target
coordinate is fixed at $X_k=t$.  Let
\[
  Z_\lambda=W_{z_\lambda}(X_I-m_{\lambda,I}).
\]
The conditional bulk mean shift is $z_\lambda\sigma_\lambda^{-1}\Delta_t$;
therefore
\[
  \mathbb{E}_{\mu_\lambda^t}Z_\lambda
  =\langle z_\lambda,z_\lambda\sigma_\lambda^{-1}\Delta_t
    \rangle_{H_{C_\lambda}}
  =\rho_\lambda\Delta_t.
\]
The conditional bulk covariance is
$C_\lambda-z_\lambda\sigma_\lambda^{-1}z_\lambda^*$, so
\[
  \operatorname{Var}_{\mu_\lambda^t}(Z_\lambda)
  =\|z_\lambda\|_{H_{C_\lambda}}^2
   -\sigma_\lambda^{-1}\|z_\lambda\|_{H_{C_\lambda}}^4
  =\sigma_\lambda\rho_\lambda(1-\rho_\lambda).
\]
Since
$X_k-\mu_{k|I}(X_I)=\Delta_t-Z_\lambda$ under $\mu_\lambda^t$,
\begin{align*}
  \mathbb{E}_{\mu_\lambda^t}
  \left[(X_k-\mu_{k|I}(X_I))^2\right]
  &=\sigma_\lambda\rho_\lambda(1-\rho_\lambda)
    +\Delta_t^2(1-\rho_\lambda)^2.
\end{align*}
Dividing by
$2S_{k,\lambda}=2\sigma_\lambda(1-\rho_\lambda)$ gives
\eqref{eq:residual_t_algebra}.  Adding
\eqref{eq:bulk_KL_t_algebra} and \eqref{eq:residual_t_algebra} cancels the
$\rho_\lambda/2$ terms and also combines the two signal factors:
\[
  \frac12\frac{\Delta_t^2}{\sigma_\lambda}\rho_\lambda
  +\frac12\frac{\Delta_t^2}{\sigma_\lambda}(1-\rho_\lambda)
  =\frac{\Delta_t^2}{2\sigma_\lambda}.
\]
This proves \eqref{eq:finite_part_exact_rank_one}.
\end{proof}

\begin{lemma}[Weak scheme-independence]
\label{lem:weak_scheme_independence}
For any two evidence levels $t_1,t_2$ for which the finite part is defined,
\begin{equation}
    \mathfrak{J}_\lambda(t_2)-\mathfrak{J}_\lambda(t_1)
    =
    \frac{(t_2-m_{\lambda,k})^2-(t_1-m_{\lambda,k})^2}
    {2\Sigma_{\lambda,kk}}.
    \label{eq:weak_scheme_independence_app}
\end{equation}
Replacing the Gaussian Dirac mollifier in Proposition~\ref{prop:mollification_limit} by
any centered approximate identity with finite second moment changes the extracted
finite part only by an additive term depending on the kernel, $\lambda$, and
$S_{k,\lambda}$, not on $t$.
\end{lemma}

\begin{proof}
The exact identity \eqref{eq:finite_part_exact_rank_one} shows that the only
$t$-dependent term is $\Delta_t^2/(2\sigma_\lambda)$.  The logarithmic term
$-\frac12\log(1-\rho_\lambda)$ depends on $C_\lambda$ and $\lambda$, but not on
$t$, and therefore cancels in the difference.  The same chain-rule computation
as in Proposition~\ref{prop:mollification_limit} shows that changing the centered
mollifier changes only the one-dimensional normalization and centered second
moment contribution, which are independent of the evidence level $t$.
\end{proof}

\begin{proof}[Proof of Theorem~\ref{thm:finite_part_regularization}]
Apply Lemma~\ref{lem:kl_finite_part_algebra} with $t=x_k^*$ and
$\Delta_\lambda=x_k^*-m_{\lambda,k}$.  Assumption~\ref{assum:gaussian_clamp}
gives $m_{\lambda,k}-x_k^*=b_k/\lambda$, hence
$\Delta_\lambda^2=b_k^2/\lambda^2$.  Since
$\sigma_\lambda=\varepsilon_0/(2\lambda^2)$,
\[
  \frac{\Delta_\lambda^2}{2\sigma_\lambda}
  =\frac{b_k^2}{\varepsilon_0}.
\]
By Lemma~\ref{lem:feldman_hajek}, $\rho_\lambda=O(\lambda^{-4})$, and therefore
\[
  -\frac12\log(1-\rho_\lambda)=O(\lambda^{-4}).
\]
The exact finite-part identity \eqref{eq:finite_part_exact_rank_one} gives
\[
  \mathfrak{J}_\lambda(\mu_{\mathrm{obs}};\mu_\lambda)
  =
  \frac{b_k^2}{\varepsilon_0}+O(\lambda^{-4}).
\]
\end{proof}

If $a_{Ik}=0$ and $x_k^*=m_{\lambda,k}$, then
$z_\lambda=0$, the bulk marginals coincide, and the conditional residual term
in Definition~\ref{def:unified_action} is zero.  Hence
$\mathfrak{J}_\lambda(\mu_{\mathrm{obs}};\mu_\lambda)=0$ exactly for every
$\lambda>0$.

\section{Admissibility Boundary and Elliptic Benchmark}
\label{app:admissibility_boundary}

The finite-energy assumption is a boundary condition for the theory, not a
technical artifact.  The off-diagonal channel must land in the Cameron--Martin
range of the background covariance.  If it does not, the variational subtraction
need not be finite and the $O(\lambda^{-6})$ rank-one residual estimate has no
meaning.

The elliptic example of \Cref{sec:spde} makes this boundary concrete.  For
$L=-d^2/dx^2+\kappa^2$ on $L^2(0,1)$ with Dirichlet boundary conditions and
$Q_0=L^{-1}$, the Cameron--Martin space is
\[
  H_{Q_0}=H_0^1(0,1),
  \qquad
  \|h\|_{Q_0}^2
  =
  \int_0^1\bigl(|h'(x)|^2+\kappa^2|h(x)|^2\bigr)\,dx .
\]
Smooth spatially averaged observations are admissible when their representing
vectors lie in $H_0^1(0,1)$; for example, any fixed
$g\in C_c^\infty(0,1)$ has finite energy.  A point-evaluation channel is the
opposite limit.  If $g_\varepsilon$ is a standard compactly supported mollifier
concentrating at $x_0\in(0,1)$, then
\[
  \|g_\varepsilon\|_{Q_0}^2
  =
  \int_0^1\bigl(|g_\varepsilon'(x)|^2+\kappa^2|g_\varepsilon(x)|^2\bigr)\,dx
  \longrightarrow \infty
  \qquad(\varepsilon\downarrow0).
\]
Thus each fixed smoothed sensor can satisfy the finite-energy condition, but
the Dirac point-evaluation limit does not.  The shorting theorem is therefore a
theory of admissible Cameron--Martin channels, not of unsmoothed point sensors
outside the covariance-energy domain.
